\section{Desarrollo de la implementación}

Esta sección describe cómo se construyó el simulador de tráfico
unidimensional con paralelismo OpenMP, explicando tanto las decisiones
algorítmicas como las consecuencias computacionales de cada una. La
premisa central fue mantener una única base de código para la lógica de
simulación y variar exclusivamente el mecanismo de ejecución y el nivel
de optimización aplicado, de forma que cualquier diferencia de rendimiento
sea atribuible al grado de paralelismo o a la estrategia de optimización
y no a diferencias en el algoritmo.

% ============================================================
\subsection{Modelo del autómata celular de tráfico}
% ============================================================

El sistema modela una carretera circular de $N$ celdas con condiciones
de contorno periódicas. Cada celda $i \in \{0, 1, \ldots, N-1\}$ puede
estar vacía ($R = 0$) o contener un carro ($R = 1$). La velocidad máxima
es $v_{\max} = 1$, por lo que en cada paso de tiempo un carro puede
avanzar a la celda siguiente si esta está libre, o permanecer bloqueado si
está ocupada. La regla de actualización se define formalmente como

\begin{equation}
    R(t+1,\, i) =
    \bigl[R(t,i) \wedge R(t,\,(i+1)\bmod N)\bigr]
    \;\vee\;
    \bigl[R(t,\,(i-1)\bmod N) \wedge \neg R(t,i)\bigr].
    \label{eq:regla}
\end{equation}

El primer término captura el caso en que el carro en $i$ no puede avanzar
porque la celda siguiente está ocupada. El segundo captura el caso en que
un carro en $i-1$ avanza a la celda $i$, que estaba libre. La
condición $\neg R(t,i)$ en el segundo término es necesaria para evitar
que una celda reciba dos carros simultáneamente.

En la implementación base, los valores booleanos se almacenan como enteros
con signo de 32~bits (\texttt{int}). La regla~\eqref{eq:regla} se traduce
directamente a operaciones de bit:

\begin{equation}
    \texttt{next}[i] =
    (\textit{here} \;\&\; \textit{ahead})
    \;\mid\;
    (\textit{behind} \;\&\; \mathtt{\sim}\textit{here} \;\&\; 1),
\end{equation}

donde \texttt{\&\,1} enmascara el complemento para suprimir la extensión
de signo que produce \texttt{\textasciitilde} sobre enteros de 32~bits con
signo; sin esa máscara, \texttt{\textasciitilde 0} produce
\texttt{0xFFFFFFFF} en lugar de~1, lo que corrompería el conteo de carros.

Todas las lecturas de la generación $t$ provienen del arreglo
\texttt{cells} y todas las escrituras de la generación $t+1$ van al
arreglo \texttt{next\_cells}. Como ambos arreglos son distintos y cada
posición $i$ de la nueva generación depende exclusivamente de posiciones
de la generación anterior, todas las iteraciones del bucle son
independientes entre sí y no existe ninguna dependencia de datos.

% ============================================================
\subsection{Velocidad media y estado estacionario}
% ============================================================

La velocidad media en el paso $t$ se define como la fracción de carros
que se desplazaron:

\begin{equation}
    v(t) = \frac{1}{N_c}
    \sum_{i=0}^{N-1}
    \mathbf{1}\bigl[R(t,i) = 1 \;\wedge\; R(t+1,i) = 0\bigr],
    \label{eq:velocidad}
\end{equation}

donde $N_c$ es el número total de carros en la carretera, cantidad
conservada: la regla~\eqref{eq:regla} nunca crea ni destruye carros. Esta
conservación se verifica explícitamente en los tests de corrección.

Partiendo de una distribución inicial aleatoria, el sistema atraviesa un
transitorio durante el cual $v(t)$ fluctúa mientras los carros se
reorganizan. Después de suficientes pasos, $v(t)$ converge a la velocidad
asintótica $v_\infty$, que es la cantidad de interés físico. El benchmark
reporta el promedio sobre la fase de medición:

\begin{equation}
    v_\infty \approx \bar{v} =
    \frac{1}{T_m}
    \sum_{t = T_w + 1}^{T_w + T_m} v(t),
    \label{eq:vasint}
\end{equation}

donde $T_w$ es el número de pasos de calentamiento y $T_m$ es el número
de pasos de medición.

% ============================================================
\subsection{Criterio de convergencia del calentamiento}
% ============================================================

La duración del transitorio no es fija: depende tanto de la densidad como
de $N$. Un calentamiento con número de pasos fijo puede ser insuficiente
para densidades cercanas al punto de transición de fase ($\rho \approx
0.5$), donde el transitorio es máximo, o innecesariamente largo para
densidades extremas donde el sistema converge en pocas decenas de pasos.

Para detectar la convergencia de forma adaptativa, se compara la velocidad
media en dos ventanas temporales contiguas de $W$ pasos. La condición de
parada es

\begin{equation}
    \left|
        \bar{v}_{[t-W,\, t]}
        - \bar{v}_{[t-2W,\, t-W]}
    \right| < \varepsilon,
    \label{eq:criterio}
\end{equation}

con $W = 50$, $\varepsilon = 10^{-4}$ y un mínimo de 200~pasos antes de
evaluar el criterio para evitar falsos positivos durante la fase inicial.
El parámetro \texttt{max\_warmup\_steps}, pasado como argumento al
binario, actúa como techo; si la condición~\eqref{eq:criterio} no se
alcanza antes de este límite, el calentamiento termina igualmente. En el
script de benchmark este techo se fija como $\max(N / 10,\; 2000)$, valor
que garantiza al menos una revolución completa de información alrededor de
la carretera. El número de pasos realmente ejecutados se registra en
\texttt{stderr} como diagnóstico, sin contaminar el \texttt{stdout} que
leen los scripts.

% ============================================================
\subsection{Arquitectura de datos}
% ============================================================

La estructura central es \texttt{Road}, que encapsula el estado completo
de la simulación:

\begin{itemize}
    \item \texttt{cells}: puntero al estado de la generación actual, usado
          exclusivamente para lectura dentro de un paso.
    \item \texttt{next\_cells}: puntero al buffer de la generación
          siguiente, usado exclusivamente para escritura.
    \item \texttt{length}: número de celdas $N$.
    \item \texttt{car\_count}: número de carros $N_c$, invariante durante
          toda la simulación.
\end{itemize}

Ambos buffers se reservan en \texttt{road\_create} con una única llamada a
\texttt{malloc} por buffer. No se realiza ninguna asignación de memoria
dentro del bucle de simulación, lo que elimina el ruido del asignador
dinámico de las mediciones de tiempo. Al finalizar cada paso, los
punteros \texttt{cells} y \texttt{next\_cells} se intercambian mediante
una permutación de punteros de costo $O(1)$, sin copiar datos.

\begin{figure}[H]
\centering
\fbox{%
\begin{minipage}{0.92\textwidth}
\centering
\textbf{Flujo de un paso de simulación}\\[4pt]
Leer \texttt{cells}
$\rightarrow$ Escribir \texttt{next\_cells} (regla~\eqref{eq:regla})
$\rightarrow$ Contar carros desplazados (ecuación~\eqref{eq:velocidad})
$\rightarrow$ Intercambiar punteros\\[2pt]
Si fase de calentamiento: evaluar criterio~\eqref{eq:criterio};
si no: acumular $v(t)$
\end{minipage}}
\caption{Ciclo compartido por todas las implementaciones.}
\label{fig:flujo_paso}
\end{figure}

El estado inicial de \texttt{cells} se genera con \texttt{rand} en
\texttt{road\_create}, sembrando el generador con \texttt{time(NULL)}.
La implementación no controla la semilla explícitamente porque la magnitud
de interés ($v_\infty$) es estadísticamente robusta respecto a la
distribución inicial, como corroboran las varianzas de
\texttt{avg\_velocity} entre repeticiones en el CSV de resultados.

% ============================================================
\subsection{Implementación secuencial}
% ============================================================

El binario \texttt{traffic\_seq} se compila sin el flag \texttt{-fopenmp}.
Las directivas \texttt{\#pragma omp} presentes en \texttt{simulation.c}
son ignoradas por el compilador conforme a ISO/IEC~9899:2011,
sección~6.10.6, que autoriza al preprocesador a descartar cualquier pragma
no reconocido. Como consecuencia, el binario no enlaza \texttt{libgomp},
no crea ningún equipo de hilos en tiempo de ejecución y el bucle de
actualización se ejecuta como un bucle \texttt{for} ordinario. Esta
propiedad puede verificarse inspeccionando las dependencias del binario con
\texttt{ldd bin/traffic\_seq}, que no lista \texttt{libgomp.so}.

La medición de tiempo utiliza
\texttt{clock\_gettime(CLOCK\_MONOTONIC)} de \texttt{time.h}, accesible
bajo \texttt{\_POSIX\_C\_SOURCE=200809L} con \texttt{-std=c11}. Esta macro
de prueba de características se define en \texttt{CFLAGS} del
\texttt{Makefile} en lugar de en los encabezados, siguiendo la práctica
estándar de que las macros de prueba pertenecen al entorno de compilación
y no al código fuente. El reloj \texttt{CLOCK\_MONOTONIC} es inmune a
ajustes del reloj del sistema (\textsc{ntp}, \texttt{adjtime}) y garantiza
mediciones de tiempo transcurrido fiables. El cronómetro se activa
inmediatamente antes de la fase de medición y se detiene al finalizar,
excluyendo el calentamiento del tiempo reportado.

La interfaz de línea de comandos es:

\begin{center}
\texttt{traffic\_seq} $N$ $\rho$ \texttt{max\_warmup} $T_m$
\end{center}

La salida estándar contiene exclusivamente dos valores en una sola línea:
el tiempo de pared en milisegundos y la velocidad media $\bar{v}$,
separados por un espacio. Esta separación entre tiempo y velocidad permite
al script de benchmark detectar automáticamente ejecuciones que no
alcanzaron el estado estacionario: si $\bar{v}$ presenta alta varianza
entre repeticiones de los mismos parámetros, el techo de calentamiento fue
insuficiente.

% ============================================================
\subsection{Implementación paralela con OpenMP}
% ============================================================

El binario \texttt{traffic\_omp} se compila con \texttt{-fopenmp}, que
activa el preprocesador de pragmas OpenMP, define la macro
\texttt{\_OPENMP} y enlaza \texttt{libgomp}. El número de hilos no se fija
en tiempo de compilación sino que se lee como argumento de línea de
comandos y se comunica al runtime mediante \texttt{omp\_set\_num\_threads},
declarada en \texttt{omp.h}. Esto permite que un único binario sirva para
todos los experimentos de escalado fuerte sin necesidad de recompilación.

La interfaz amplía la del binario secuencial con un argumento adicional:

\begin{center}
\texttt{traffic\_omp} $N$ $\rho$ \texttt{max\_warmup} $T_m$ $p$
\end{center}

donde $p$ es el número de hilos solicitado. El valor efectivo se confirma
con \texttt{omp\_get\_max\_threads()} y se registra en \texttt{stderr}.

Dentro de \texttt{simulation.c} existen dos regiones paralelas
independientes por paso de tiempo. La primera, en
\texttt{compute\_next\_generation}, aplica la regla~\eqref{eq:regla} a
cada celda:

\begin{center}
\begin{verbatim}
#pragma omp parallel for schedule(static)
for (int i = 0; i < length; i++) { ... }
\end{verbatim}
\end{center}

La cláusula \texttt{schedule(static)} asigna bloques contiguos de celdas
a cada hilo con partición determinista en tiempo de compilación. Esta
elección es apropiada porque el trabajo por celda es uniforme: la
regla~\eqref{eq:regla} implica exactamente tres lecturas, dos operaciones
lógicas y una escritura para cada posición, independientemente del estado
de la celda. Con trabajo uniforme, \texttt{schedule(static)} tiene menor
overhead de planificación que \texttt{dynamic} o \texttt{guided}, y
produce un patrón de acceso a memoria favorable para el prefetcher al
asignar rangos contiguos a cada hilo.

La independencia de las iteraciones es estructural: todas las lecturas
provienen de \texttt{current[]} y todas las escrituras van a
\texttt{next[]}, que son punteros a arreglos distintos. No existe ninguna
dependencia de datos entre iteraciones, lo que garantiza ausencia de
condiciones de carrera sin necesidad de secciones críticas.

La segunda región paralela, en \texttt{count\_departed\_cars}, cuenta los
carros que se desplazaron para calcular $v(t)$:

\begin{center}
\begin{verbatim}
#pragma omp parallel for reduction(+:departed) schedule(static)
for (int i = 0; i < length; i++)
    departed += (old_cells[i] == 1 && new_cells[i] == 0) ? 1 : 0;
\end{verbatim}
\end{center}

La cláusula \texttt{reduction(+:departed)} instruye a OpenMP a mantener
una copia privada del acumulador en cada hilo y combinarlas con suma al
finalizar la región paralela. Sin esta cláusula, los incrementos
concurrentes sobre una variable compartida constituirían una condición de
carrera. Entre las dos regiones paralelas no se requiere sincronización
explícita: la barrera implícita al final de la primera directiva garantiza
que todos los hilos han completado sus escrituras en \texttt{next\_cells}
antes de que cualquier hilo comience \texttt{count\_departed\_cars}.

% ============================================================
\subsection{Optimización por compilador}
% ============================================================

Para cuantificar la contribución del compilador al rendimiento, se
construyen dos variantes adicionales — \texttt{traffic\_seq\_opt} y
\texttt{traffic\_omp\_opt} — a partir de los \emph{mismos} archivos
fuente que las versiones base, variando exclusivamente el conjunto de
flags de compilación. Esta metodología garantiza que cualquier diferencia
de rendimiento sea atribuible íntegramente al compilador y no a cambios
algorítmicos.

El conjunto de flags optimizadas se seleccionó específicamente para el
perfil de este kernel. El loop crítico de \texttt{compute\_next\_generation}
realiza tres cargas de memoria, dos operaciones bitwise y una escritura por
celda, lo que produce una intensidad aritmética de aproximadamente $0.4$
operaciones por byte de memoria transferida. Con este perfil predominantemente
\emph{memory-bound}, las flags con mayor impacto son las que amplían el
ancho de los vectores SIMD y reducen el overhead de llamada entre
unidades de traducción:

\begin{itemize}
    \item \texttt{-O3}: habilita vectorización automática
          (\texttt{-ftree-vectorize}), transformaciones de bucle e
          inlining agresivo.

    \item \texttt{-march=native}: emite instrucciones AVX2 (256~bits) o
          AVX-512 (512~bits) según el procesador donde se compila.
          Con datos \texttt{int} de 32~bits, un registro AVX2 procesa 8
          celdas por instrucción; con \texttt{uint8\_t} (variante de
          optimización de memoria, descrita más adelante), procesa~32.

    \item \texttt{-funroll-loops}: reduce el overhead del contador del
          bucle; beneficia especialmente los tamaños L1-bound ($N \leq
          64\,000$).

    \item \texttt{-flto} con \texttt{-fno-semantic-interposition}:
          Link-Time Optimization permite al compilador inlinar
          \texttt{compute\_next\_generation} y
          \texttt{count\_departed\_cars} directamente en
          \texttt{simulation\_step}, eliminando el overhead de llamada
          en cada uno de los $T_m = 1000$ pasos por medición. Sin
          \texttt{-fno-semantic-interposition}, GCC emite llamadas
          indirectas a funciones externas para soportar
          \texttt{LD\_PRELOAD}, lo que impide el inlining entre unidades
          de traducción.

    \item \texttt{-falign-loops=32}: alinea el inicio de cada bucle a un
          boundary de 32~bytes, evitando que el cuerpo del bucle cruce
          dos líneas de instrucciones de AVX2.
\end{itemize}

El flag \texttt{-ffast-math} se omite deliberadamente: el bucle crítico
no usa punto flotante y su inclusión podría hacer que \texttt{avg\_velocity}
no sea reproducible entre compiladores, complicando la validación cruzada
del autómata.

Las versiones base (\texttt{traffic\_seq}, \texttt{traffic\_omp}) se
compilan con \texttt{-O0} para constituir una línea base sin ninguna
optimización, lo que permite aislar el efecto del compilador con máxima
claridad.

% ============================================================
\subsection{Optimización de memoria}
% ============================================================

El análisis de profiling reveló que el cuello de botella del kernel no es
la lógica de cómputo, el IPC medido es de aproximadamente $3.7$, cercano
al límite teórico del superescalar de 4~vías, sino el ancho de banda de
memoria. La causa estructural es el tipo de dato: cada celda ocupa
\texttt{sizeof(int) = 4}~bytes pero almacena únicamente los valores $0$ o
$1$, dejando 3~bytes sin utilizar por cada celda. Una línea de caché de
64~bytes transporta 16~celdas con la representación \texttt{int}, cuando
podría transportar 64~con representación \texttt{uint8\_t}. Para
$N = 2\,000\,000$, el working set con \texttt{int} es de 16~MB y
desborda la L3; con \texttt{uint8\_t} se reduce a 4~MB, que cabe
completamente en caché.

Para cuantificar este efecto se implementaron dos variantes adicionales,
\texttt{traffic\_seq\_mem} y \texttt{traffic\_omp\_mem\_opt}, que
aplican las siguientes tres optimizaciones de forma conjunta.

\subsubsection{MEM-1: Reducción del tipo de celda a \texttt{uint8\_t}}

La estructura \texttt{RoadOpt} reemplaza los punteros \texttt{int~*} por
\texttt{uint8\_t~*}. La regla~\eqref{eq:regla} en código de bits se
preserva sin cambios lógicos, pero se añade un cast explícito al
complemento:

\begin{equation}
    \texttt{next}[i] =
    (h \;\&\; a) \;\mid\;
    (b \;\&\; \mathtt{(uint8\_t)(\sim h)} \;\&\; \mathtt{1u}),
\end{equation}

donde el cast \texttt{(uint8\_t)} evita que la extensión de signo de
\texttt{\textasciitilde h} sobre 32~bits corrompa el resultado al volver a
asignarlo en un campo de 8~bits. El working set se reduce en un factor de
4, con el impacto sobre la jerarquía de caché descrito en la
Tabla~\ref{tab:ws}.

\begin{table}[H]
\centering
\caption{Working set con \texttt{int} vs \texttt{uint8\_t} para dos buffers.}
\label{tab:ws}
\begin{tabular}{rrrrl}
\hline
$N$ & \texttt{int} (bytes) & \texttt{uint8\_t} (bytes)
    & Nivel \texttt{int} & Nivel \texttt{uint8\_t} \\
\hline
8\,000       & 64\,KB    & 16\,KB  & L1      & L1 \\
64\,000      & 512\,KB   & 128\,KB & L2      & L1 \\
500\,000     & 4\,MB     & 1\,MB   & L3      & L2/L3 \\
2\,000\,000  & 16\,MB    & 4\,MB   & DRAM    & L3 \\
5\,000\,000  & 40\,MB    & 10\,MB  & DRAM    & DRAM \\
10\,000\,000 & 80\,MB    & 20\,MB  & DRAM    & DRAM \\
\hline
\end{tabular}
\end{table}

\subsubsection{MEM-2: Fusión de pasadas y eliminación del módulo}

La implementación base realiza dos pasadas completas por paso de
simulación: \texttt{compute\_next\_generation} recorre los $N$ elementos
para producir \texttt{next[]}, y \texttt{count\_departed\_cars} los
recorre de nuevo para contar los carros desplazados. Para tamaños mayores
que la L3, los datos escritos en la primera pasada son desalojados de caché
antes de que comience la segunda, duplicando el tráfico hacia DRAM. La
fusión de ambas pasadas en \texttt{compute\_next\_and\_count} lee y escribe
cada celda una sola vez por paso:

\begin{center}
\begin{verbatim}
/* bucle interno — sin módulo, vectorizable */
#pragma omp parallel for schedule(static) reduction(+:departed)
for (int i = 1; i < length - 1; i++) {
    uint8_t b = current[i - 1];
    uint8_t h = current[i];
    uint8_t a = current[i + 1];
    uint8_t n = (h & a) | (b & (uint8_t)(~h) & 1u);
    next[i]   = n;
    departed += (h == 1u && n == 0u) ? 1 : 0;
}
\end{verbatim}
\end{center}

Las celdas frontera ($i = 0$ e $i = N-1$), que requieren índices con
módulo, se calculan secuencialmente fuera de la región paralela antes y
después del bucle interno. Esta técnica de \emph{boundary peeling} elimina
el operador módulo, que implica una división entera implícita, del camino
crítico y deja el bucle interno con índices $i{-}1$, $i$ e $i{+}1$
directos, lo que permite al vectorizador automático emitir cargas e
instrucciones AVX2 sin bifurcaciones de índice.

\subsubsection{MEM-3: Alineación de buffers a línea de caché}

Los buffers se reservan con \texttt{posix\_memalign(64, \ldots)} en lugar
de \texttt{malloc}. Esto garantiza que el primer byte de cada arreglo
coincida con el inicio de una cacheline de 64~bytes, con dos efectos:
primero, los accesos del primer y último hilo OMP al extremo de su rango
asignado no producen \emph{false sharing} con la línea de caché inicial o
final del buffer; segundo, el hardware stream prefetcher inicia el
seguimiento del patrón de acceso desde una dirección base alineada, lo
que reduce la latencia de la primera línea traída del siguiente nivel de
la jerarquía. La memoria devuelta por \texttt{posix\_memalign} es
compatible con \texttt{free}, por lo que \texttt{road\_mem\_opt\_destroy}
no requiere cambios respecto a \texttt{road\_destroy}.

% ============================================================
\subsection{Matriz de variantes de compilación}
% ============================================================

Las cuatro estrategias de optimización anteriores se combinan en ocho
binarios que cubren todos los ejes del espacio de diseño:

\begin{table}[H]
\centering
\caption{Binarios generados por \texttt{make all} y sus características.}
\label{tab:binarios}
\begin{tabular}{llccc}
\hline
Binario & \texttt{impl} en CSV & OMP & OPT flags & Mem opt \\
\hline
\texttt{traffic\_seq}         & \texttt{traffic\_seq}         & No  & No  & No \\
\texttt{traffic\_seq\_opt}    & \texttt{traffic\_seq\_opt}    & No  & Sí  & No \\
\texttt{traffic\_omp}         & \texttt{traffic\_omp}         & Sí  & No  & No \\
\texttt{traffic\_omp\_opt}    & \texttt{traffic\_omp\_opt}    & Sí  & Sí  & No \\
\texttt{traffic\_seq\_mem}    & \texttt{traffic\_seq\_mem}    & No  & No  & Sí \\
\texttt{traffic\_seq\_mem\_opt} & \texttt{traffic\_seq\_mem\_opt} & No & Sí & Sí \\
\texttt{traffic\_omp\_mem}    & \texttt{traffic\_omp\_mem}    & Sí  & No  & Sí \\
\texttt{traffic\_omp\_mem\_opt} & \texttt{traffic\_omp\_mem\_opt} & Sí & Sí & Sí \\
\hline
\end{tabular}
\end{table}

% ============================================================
\subsection{Métricas de rendimiento}
% ============================================================

El script de benchmark registra el wall time $T(p)$ de la fase de
medición para cada configuración $(N, \rho, p, \text{impl})$. A partir
del CSV generado, las métricas de escalado fuerte se calculan como:

\begin{equation}
    S(p) = \frac{T_{\text{seq}}}{T_{\text{omp}}(p)},
    \qquad
    E(p) = \frac{S(p)}{p} \times 100\%.
    \label{eq:speedup}
\end{equation}

La eficiencia $E(p)$ cuantifica la fracción de capacidad paralela
efectivamente utilizada. Para este autómata, cuyo kernel tiene una
intensidad aritmética de aproximadamente $0.4$~operaciones por byte de
memoria transferida, se espera que $E(p)$ decrezca con $N$ pequeño (donde
el overhead de creación del equipo de hilos domina) y con $N$ muy grande
(donde el ancho de banda de memoria es el cuello de botella y agregar
hilos no acelera la transferencia de datos). La optimización de memoria
desplaza hacia arriba el umbral de $N$ a partir del cual la eficiencia
paralela se degrada, al reducir el tráfico por celda en un factor de~4.

% ============================================================
\subsection{Diseño del script de pruebas}
% ============================================================

El script \texttt{bench\_traffic.sh} actúa como harness de medición: su
única responsabilidad es invocar los binarios con los parámetros correctos,
capturar su salida y persistirla en un CSV. 

% ============================================================
\subsection{Matriz de configuraciones}
% ============================================================

El harness define dos experimentos que cubren las preguntas de rendimiento
centrales del trabajo.

El \emph{experimento de escalado} mide todas las combinaciones de los seis
valores de $N$ con los seis conteos de hilos $\{0, 2, 4, 6, 8, 12\}$ a
densidad fija $\rho = 0.50$, para los ocho binarios descritos en la
Tabla~\ref{tab:binarios}. El valor $p = 0$ indica los binarios secuenciales;
cualquier $p > 0$ invoca los binarios OpenMP con ese número de hilos. Los
valores de $N$ se eligen para cubrir cada nivel de la jerarquía de
memoria con la representación base (\texttt{int}); la
Tabla~\ref{tab:ws} muestra cómo la representación \texttt{uint8\_t} desplaza
el punto de operación de varios tamaños hacia niveles superiores de caché:

\begin{center}
\begin{tabular}{rrl}
\hline
$N$ & Working set (\texttt{int}) & Nivel esperado \\
\hline
8\,000       & 64\,KB  & L1 \\
64\,000      & 512\,KB & L2 \\
500\,000     & 4\,MB   & L3 superior \\
2\,000\,000  & 16\,MB  & Frontera L3/DRAM \\
5\,000\,000  & 40\,MB  & DRAM (referencia de escalado fuerte) \\
10\,000\,000 & 80\,MB  & DRAM saturada \\
\hline
\end{tabular}
\end{center}

El \emph{experimento de densidad} fija $N = 20\,000\,000$ (160~MB con
\texttt{int}; 40~MB con \texttt{uint8\_t}) y mide los hilos
$\{0, 4, 8, 12\}$ sobre los cinco valores de densidad
$\{0.10,\, 0.30,\, 0.50,\, 0.70,\, 0.90\}$. Este valor de $N$ es
deliberadamente exterior al conjunto $N\_\mathrm{VALUES}$ del experimento
de escalado para que las dos series de datos sean independientes en el
CSV. Para los hilos secuenciales ($p = 0$) se miden las variantes
\texttt{traffic\_seq}, \texttt{traffic\_seq\_opt} y
\texttt{traffic\_seq\_mem\_opt}; para los hilos paralelos, únicamente
\texttt{traffic\_omp\_mem\_opt}, que representa la configuración más
eficiente y cuyo comportamiento ante distintas densidades es el de mayor
interés práctico.

% ============================================================
\subsection{Techo de calentamiento adaptativo}
% ============================================================

El techo de pasos de calentamiento que el harness pasa a cada binario se
calcula como

\begin{equation}
    T_{w,\max} = \max\!\left(\frac{N}{10},\; 2000\right).
    \label{eq:techo}
\end{equation}

El factor $N/10$ garantiza que el techo sea al menos proporcional al
tiempo que tarda una perturbación en recorrer la carretera completa, del
orden de $N$ pasos, con un margen de seguridad de un orden de magnitud.
El piso de 2000 evita techos irrazonablemente bajos para $N$ pequeños.
El binario puede terminar el calentamiento antes de alcanzar este techo si
el criterio~\eqref{eq:criterio} se satisface, lo que ocurre en la
práctica para densidades extremas ($\rho \leq 0.1$ o $\rho \geq 0.9$) en
pocas centenas de pasos.


% ============================================================
\subsection{Formato del CSV de resultados}
% ============================================================

El archivo \texttt{data.csv} tiene siete columnas:

\begin{center}
\begin{tabular}{ll}
\hline
Columna & Descripción \\
\hline
\texttt{impl}          & Binario usado (véase Tabla~\ref{tab:binarios}) \\
\texttt{threads}       & Número de hilos (0 para variantes secuenciales) \\
\texttt{road\_length}  & Tamaño de la carretera $N$ \\
\texttt{density}       & Densidad inicial de carros $\rho$ \\
\texttt{repetition}    & Índice de repetición (1 a 10) \\
\texttt{wall\_time\_ms} & Tiempo de pared de la fase de medición en ms \\
\texttt{avg\_velocity} & Velocidad media $\bar{v}$ sobre los $T_m$ pasos \\
\hline
\end{tabular}
\end{center}

La columna \texttt{avg\_velocity} cumple una función de diagnóstico: para
una combinación $(N, \rho, \text{impl})$ dada, las 10 repeticiones deben
producir valores de $\bar{v}$ con varianza inferior a $10^{-3}$ si el
sistema alcanzó el estado estacionario antes de iniciar la medición.
Adicionalmente, para los mismos parámetros $(N, \rho)$, todos los binarios
deben coincidir en $\bar{v}$ independientemente del número de hilos o del
nivel de optimización: todos implementan el mismo algoritmo determinista
sobre la misma condición inicial (mismo \texttt{road\_length} y misma
semilla de densidad), por lo que una discrepancia sistemática indicaría
una condición de carrera en la versión OpenMP o una corrupción de datos
introducida por las optimizaciones de memoria.